\section{Выводы и обсуждение результатов}

В ходе выполнения лабораторной работы были изучены основные свойства голограмм точечного источника и объёмного предмета. В части А работы мы определили цену деления предметной шкалы двумя методами — дифракции Фраунгофера и с помощью собирающей линзы, получив хорошее совпадение результатов с относительной погрешностью 1.8\%. 

Были экспериментально определены следующие ключевые величины:
\begin{itemize}
    \item Цена деления шкалы: $D_{\text{дифр}} = 0.11$ мм, $D_{\text{линза}} = 0.108$ мм
    \item Расстояние от голограммы до точечного источника: $d = (6.9 \pm 3.5)$ мм
    \item Расстояние до действительного изображения: $d_{\text{действ}} = (35.2 \pm 2.0)$ мм
    \item Расстояние до мнимого изображения: $d_{\text{мним}} = (27.7 \pm 2.0)$ мм
    \item Фокусное расстояние голографической линзы: $F = 4.27$ см
\end{itemize}

В результате исследования фокусирующих свойств голограммы было установлено, что левое пятно большего размера соответствует мнимому изображению, а правое меньшего размера — действительному. При смещении голограммы перпендикулярно оптической оси было обнаружено, что расстояния до изображений остаются неизменными, а сами изображения смещаются, сохраняя свои характеристики, что принципиально отличает голограмму от обычной линзы.

В части Б работы были исследованы характеристики голограммы объёмного предмета. Мы наблюдали эффект параллакса, подтверждающий трёхмерный характер голографического изображения, и определили угол падения опорной волны при записи голограммы — около 50° относительно нормали к фотопластинке. Было подтверждено фундаментальное свойство голограмм — способность восстанавливать изображение всего объекта по части голограммы.

Методом параллакса было измерено расстояние между элементами объекта (линейкой и вертикальным стержнем): $d = 1.75$ мм в масштабе голографического изображения. При исследовании действительного изображения обнаружена инверсия глубины по сравнению с мнимым изображением, а при повороте голограммы на 180° наблюдалась полная инверсия восстановленного волнового фронта.

Таким образом, экспериментально подтверждены основные теоретические положения голографии: голограмма регистрирует полную информацию о волновом фронте объекта, включая амплитуду и фазу; может формировать как мнимое, так и действительное изображение; сохраняет все пространственные соотношения между элементами объекта; обладает объёмностью и демонстрирует эффект параллакса; каждый участок голограммы содержит информацию обо всём объекте. Полученные результаты свидетельствуют о полном выполнении цели работы.
